% Algebra 2 Lesson 3-4 Lesson Quiz, short form (built from the TE "3-4 Lesson Quiz: Dividing Polynomials" block, a2_3-4_TE.tex).
% The short form keeps the Savvas item numbers 2, 3, 5. Savvas items 1 and 4 are left off (see the key).
\documentclass[11pt]{article}
\usepackage[margin=0.7in]{geometry}
\usepackage{amsmath,amssymb}
\usepackage{enumitem}
\usepackage{tikz}
\usepackage{pgfplots}
\pgfplotsset{compat=1.18}
\usepackage{fancyhdr}
\usepackage{parskip}
\usepackage{array}
\setlist[enumerate]{leftmargin=*, itemsep=6pt}
\pagestyle{fancy}
\fancyhf{}
\fancyfoot[L]{\small\textit{en}Vision\textsuperscript{\textregistered} Algebra 2}
\fancyfoot[C]{\small 3-4 Lesson Quiz, short form}
\fancyfoot[R]{\small\thepage}
\renewcommand{\headrulewidth}{0pt}
\newcommand{\choice}[1]{\tikz[baseline=-0.6ex]\node[draw, circle, inner sep=0.6pt, minimum size=3.3mm, font=\scriptsize\bfseries] {#1};\;}
\newcommand{\checkbox}{$\square$\;}
\newcolumntype{C}{>{\centering\arraybackslash$}p{0.9cm}<{$}}
\newcolumntype{A}{>{$}r<{$}}
\renewcommand{\arraystretch}{1.3}

\begin{document}

\noindent\textbf{Name}\; \rule{6cm}{0.4pt} \hfill Period \rule{1.2cm}{0.4pt} \quad Date \rule{2.2cm}{0.4pt}

\vspace{4pt}\noindent\rule{\textwidth}{0.6pt}

\vspace{2pt}\noindent{\LARGE\textbf{3-4 \; Lesson Quiz: Dividing Polynomials}} \hfill {\small short form}

\vspace{4pt}\noindent\rule{\textwidth}{0.6pt}

\vspace{6pt}
\begin{minipage}[t]{0.56\textwidth}
\textbf{2.} Morgan begins to use synthetic division to divide \mbox{$x^4+x^3-6x^2-4x+8$} by $x-2$, as shown at the right. Which of the following shows the correct values to replace the question marks, and also gives the quotient?

\vspace{4pt}
\begin{itemize}[leftmargin=1.2em, itemsep=5pt]
\item[] \choice{A} $2,\ 6,\ 0,\ -8$; \ $x^3-3x^2+4$
\item[] \choice{B} $2,\ 6,\ 0,\ -8$; \ $x^3+3x^2-4$
\item[] \choice{C} $1,\ 3,\ 0,\ -4$; \ $x^2+3x-4$
\item[] \choice{D} $1,\ 1,\ -6,\ -4$; \ $x^3+3x^2-4$
\end{itemize}
\end{minipage}\hfill
\begin{minipage}[t]{0.40\textwidth}
\vspace{4pt}
\setlength{\tabcolsep}{5pt}%
\begin{tabular}{A|CCCCC}
2 & 1 & 1 & -6 & -4 & 8 \\
  &   & ? & ?  & ?  & ? \\ \cline{2-6}
  & \rule{0pt}{18pt} & & & & \\
\end{tabular}

\vspace{6pt}
{\small Show your work: finish the synthetic division in the table.}
\end{minipage}

\vspace{10pt}\noindent\rule{\textwidth}{0.4pt}

\vspace{6pt}
\textbf{3.} Verify the Remainder Theorem if $P(x)=2x^3-x^2+4x+5$ is divided by $x+2$. Complete the sentences.

\vspace{6pt}
\begin{minipage}[t]{0.56\textwidth}
When $P(x)$ is divided by $x+2$, the remainder is

\hspace*{1em}\choice{} $-23$ \quad \choice{} $-7$ \quad \choice{} $1$ \quad \choice{} $33$

\vspace{8pt}
$P(-2)=$ \quad \choice{} $-25$ \quad \choice{} $-23$ \quad \choice{} $9$ \quad \choice{} $25$

\vspace{8pt}
Does this example verify the Remainder Theorem?

\hspace*{1em}\choice{} Yes \quad \choice{} No
\end{minipage}\hfill
\begin{minipage}[t]{0.40\textwidth}
{\small Show your work (synthetic division and $P(-2)$):}

\vspace{4pt}
\begin{tabular}{A|CCCC}
\rule{0pt}{18pt}\phantom{-2} & & & & \\
\rule{0pt}{18pt} & & & & \\ \cline{2-5}
\rule{0pt}{18pt} & & & & \\
\end{tabular}

\vspace{10pt}\rule{\linewidth}{0.4pt}
\end{minipage}

\vspace{10pt}\noindent\rule{\textwidth}{0.4pt}

\vspace{6pt}
\begin{minipage}[t]{0.56\textwidth}
\textbf{5.} If $x+2$ is a factor of $P(x)=2x^3+5x^2+5x+6$, select the two factors that express $P(x)$ as a product.

\vspace{4pt}
\begin{itemize}[leftmargin=1.2em, itemsep=5pt]
\item[] \checkbox \textbf{A.} $(x+2)$
\item[] \checkbox \textbf{B.} $(2x^2-x-3)$
\item[] \checkbox \textbf{C.} $(2x^2-x+3)$
\item[] \checkbox \textbf{D.} $(2x^2+x-3)$
\item[] \checkbox \textbf{E.} $(2x^2+x+3)$
\item[] \checkbox \textbf{F.} not a factor
\end{itemize}
\end{minipage}\hfill
\begin{minipage}[t]{0.40\textwidth}
{\small Show your work:}

\vspace{4pt}
\begin{tabular}{A|CCCC}
\rule{0pt}{18pt}\phantom{-2} & & & & \\
\rule{0pt}{18pt} & & & & \\ \cline{2-5}
\rule{0pt}{18pt} & & & & \\
\end{tabular}

\vspace{8pt}
{\small How does your work show that $x+2$ is a factor?}

\vspace{12pt}\rule{\linewidth}{0.4pt}
\end{minipage}

\vfill
\noindent\textbf{Scoring:} Question 2 and Question 5 are 10 points each. Question 3 is a bonus of up to 5 points. Total out of 20. (Question numbers are the Savvas Lesson Quiz numbers; items 1 and 4 are not on this form.)

\newpage
\noindent{\LARGE\textbf{3-4 Lesson Quiz, short form: answer key}}

\vspace{4pt}\noindent\rule{\textwidth}{0.6pt}

{\small The question numbers are the Savvas Lesson Quiz item numbers, as on the packet's ``On the quiz'' box and IXL table. Items 1 and 4 are left off (below).}

\begin{enumerate}
\item[\textbf{2.}] \textbf{B: $2,\ 6,\ 0,\ -8$; \ $x^3+3x^2-4$.} Bring down 1; $2\cdot1=2$, $1+2=3$; $2\cdot3=6$, $-6+6=0$; $2\cdot0=0$, $-4+0=-4$; $2\cdot(-4)=-8$, $8+(-8)=0$. Bottom row $1,\ 3,\ 0,\ -4 \mid 0$, so the quotient is $x^3+3x^2+0x-4$ with remainder $0$. A flips the signs of the last two quotient terms, C puts the bottom row where the products go and drops a degree, D copies the top row. Savvas quiz item 2, DOK 2, Realize skill A75. Taught by Example 2. \textbf{10 points:} 7 for B, 3 for the completed table.
\item[\textbf{3.}] \textbf{$-23$; \ $-23$; \ Yes.} Synthetic division by $x-(-2)$: bottom row $2,\ -5,\ 14 \mid -23$. Direct substitution: $P(-2)=2(-8)-4+4(-2)+5=-16-4-8+5=-23$. The remainder equals $P(-2)$, which verifies the Remainder Theorem. Distractor $25$ in the second sentence is $P(2)$, the sign error with $x+2$. Savvas quiz item 3, DOK 2, Realize skill A74. Taught by Example 5 (Methods 1 and 2 give the same number; Example 3 is cut). \textbf{Bonus 5:} all three correct 5; two correct 3; otherwise 0.
\item[\textbf{5.}] \textbf{A and E: $(x+2)(2x^2+x+3)$.} Synthetic division with $a=-2$: bottom row $2,\ 1,\ 3 \mid 0$. The remainder is $0$, so $x+2$ is a factor (Factor Theorem) and the quotient is $2x^2+x+3$.\newline Check: \mbox{$(x+2)(2x^2+x+3)=2x^3+5x^2+5x+6$}. Savvas quiz item 5, DOK 2, Realize skill A75. Taught by Example 5A. \textbf{10 points:} 4 for each correct factor and none wrong, 2 for the work showing a remainder of $0$.
\end{enumerate}

\textbf{Items left off this form} (available for a retake or Form B): Savvas item 1, use long division to divide $x^3+x^2-2x+14$ by $x+3$, build the answer from the choices $x^2$, $2x$, $-2x$, $4x$, $-8$, $4$, $10$, $-\frac{10}{x+3}$, $\frac{2}{x+3}$, $\frac{38}{x+3}$, $\frac{44}{x+3}$ (answer $x^2-2x+4+\frac{2}{x+3}$; DOK 2, skill A74; Example 1 long division, cut from the packet, though synthetic division with $a=-3$ gives the same result); Savvas item 4, the remainder when $f(x)=x^4-5x^2-6x-10$ is divided by $x-3$, choices A $-46$, B $8$, C $44$, D $224$ (answer B, $8$; DOK 2, skill F01; Example 4, rehearsed by Example 5 and IXL CHC).

\textbf{IXL replacement:} SmartScore 70 or above on D6D (item 2), CHC (bonus item 3) and VXE (item 5) replaces this quiz with 20/20.

\end{document}
